% Capítulo de integração numérica (arquivos usados neste capítulo)
\def\arqTrap{RegraTrapezios_repetida.ggb}
\def\arqSimpA{RegraSimpson13_repetida.ggb}
\def\arqSimpB{RegraSimpson38_repetida.ggb}
\def\arqRomberg{MetodoRomberg.ggb}
\def\arqGauss{QuadraturaGaussLegendre.ggb}
\chapter{Integração numérica}\label{cap:integracao}

\begin{tcolorbox}[colback=gray!4,colframe=gray!50,title={Como usar os aplicativos deste capítulo},fonttitle=\bfseries]
Ao longo do capítulo, os quadros laranja \emph{Atividade no GeoGebra} trazem um link
\textbf{[abrir o aplicativo]} que abre o arquivo correspondente no GeoGebra Classic:
\begin{itemize}[nosep]
  \item \ggbabrir{\arqTrap}{Regra dos trapézios repetida};
  \item \ggbabrir{\arqSimpA}{Regra 1/3 de Simpson repetida};
  \item \ggbabrir{\arqSimpB}{Regra 3/8 de Simpson repetida};
  \item \ggbabrir{\arqRomberg}{Método de Romberg};
  \item \ggbabrir{\arqGauss}{Quadratura de Gauss-Legendre}.
\end{itemize}
Para que os links funcionem, mantenha este PDF e a pasta \texttt{ggb} juntos, na mesma
pasta, e abra o PDF num leitor que permita abrir arquivos externos (Adobe Acrobat Reader,
Okular, SumatraPDF, Foxit). Os leitores de PDF embutidos nos navegadores costumam bloquear
esse tipo de link; nesse caso, abra os arquivos da pasta \texttt{ggb} diretamente no GeoGebra.

Os cinco aplicativos começam com a integral $\int_0^1 e^{-x^2}\,dx$, que é o exemplo
desenvolvido no texto. Em todos eles é possível digitar outra função no campo
$f(x)=$, mudar o intervalo $[a,b]$ nos controles deslizantes e usar o botão
\emph{Enquadrar o gráfico} para ajustar a janela. A janela da direita mostra as contas e tem
barras de rolagem (vertical e horizontal). O valor de referência da integral é calculado
pelo próprio GeoGebra, com precisão bem maior que a dos métodos estudados.
\end{tcolorbox}

% ==================================================================
\section{O problema da integração numérica}

O Teorema Fundamental do Cálculo resolve a integral definida quando conhecemos uma
primitiva: $\int_a^b f(x)\,dx=F(b)-F(a)$, com $F'=f$. Na prática, esse caminho falha com
frequência:
\begin{itemize}[nosep]
  \item muitas funções simples não têm primitiva expressa por funções elementares, como
        $e^{-x^2}$, $\dfrac{\operatorname{sen}x}{x}$, $\sqrt{1+\cos^2x}$ (comprimento de arco de
        uma senoide) ou $\dfrac{1}{\ln x}$;
  \item às vezes a primitiva existe, mas é tão complicada que avaliá-la custa mais (e
        acumula mais erros de arredondamento) do que aproximar a integral diretamente;
  \item muitas vezes $f$ só é conhecida numa tabela de valores medidos (velocidades,
        vazões, temperaturas) e não há fórmula para integrar.
\end{itemize}

A ideia central deste capítulo vem do Capítulo~\ref{cap:interp}: \emph{trocamos $f$ por um
polinômio interpolador $p$ e integramos $p$}, o que é sempre fácil. Como $p$ é uma
combinação dos valores $f(x_i)$ nos nós, a integral aproximada também é:
\[
\int_a^b f(x)\,dx\;\approx\;\int_a^b p(x)\,dx\;=\;\sum_{i=0}^{n}w_i\,f(x_i).
\]
Uma expressão desse tipo chama-se \emph{fórmula de quadratura}; os $x_i$ são os \emph{nós}
e os $w_i$ são os \emph{pesos}. Os pesos dependem só dos nós, não de $f$.

\begin{definicao}
Uma fórmula de quadratura tem \emph{grau de exatidão} $m$ se é exata (erro zero) para
todos os polinômios de grau menor ou igual a $m$ e não é exata para algum polinômio de
grau $m+1$. Como a fórmula e a integral são lineares, basta testar
$1,x,x^2,\dots,x^{m+1}$.
\end{definicao}

Usaremos ao longo de todo o capítulo o mesmo exemplo,
\[
I=\int_0^1 e^{-x^2}\,dx = 0{,}746824132812\dots,
\]
uma integral sem primitiva elementar (ela aparece na distribuição normal da Estatística,
por meio da função erro $\operatorname{erf}$). O valor acima foi obtido com alta precisão e
serve para medir o erro de cada método.

Dividimos os métodos em dois grupos:
\begin{enumerate}[nosep]
  \item \textbf{Fórmulas de Newton-Cotes}: nós igualmente espaçados, $x_i=a+ih$. São as regras
        dos trapézios, $1/3$ de Simpson e $3/8$ de Simpson, quase sempre usadas na forma
        \emph{repetida} (composta). A extrapolação de Richardson e o método de Romberg
        melhoram os trapézios sem novos tipos de fórmula.
  \item \textbf{Quadratura de Gauss}: os nós também são escolhidos, de modo a obter o maior
        grau de exatidão possível com poucos pontos.
\end{enumerate}

% ==================================================================
\section{Regra dos trapézios}

\subsection{Regra simples}

Interpolando $f$ nos dois extremos, $x_0=a$ e $x_1=b$, obtemos a reta
$p_1(x)=f(x_0)+\dfrac{f(x_1)-f(x_0)}{h}(x-x_0)$, com $h=b-a$. A área sob essa reta é a área
de um trapézio:
\[
\int_{x_0}^{x_1}f(x)\,dx\approx T=\frac{h}{2}\bigl[f(x_0)+f(x_1)\bigr].
\]

\begin{teorema}[Erro da regra dos trapézios simples]\label{teo:errotrap}
Se $f''$ é contínua em $[x_0,x_1]$ e $h=x_1-x_0$, existe $\eta\in(x_0,x_1)$ tal que
\[
\int_{x_0}^{x_1}f(x)\,dx-\frac{h}{2}\bigl[f(x_0)+f(x_1)\bigr]=-\frac{h^3}{12}\,f''(\eta).
\]
\end{teorema}

\begin{proof}
Pelo Teorema~\ref{teo:erro} (erro da interpolação) com dois nós,
$f(x)-p_1(x)=\dfrac{f''(\xi(x))}{2}(x-x_0)(x-x_1)$. O produto $(x-x_0)(x-x_1)$ não muda de
sinal em $[x_0,x_1]$ (é $\le0$), então o Teorema do Valor Médio para integrais permite tirar
$f''$ de dentro da integral:
\[
\int_{x_0}^{x_1}\bigl(f-p_1\bigr)\,dx=\frac{f''(\eta)}{2}\int_{x_0}^{x_1}(x-x_0)(x-x_1)\,dx
=\frac{f''(\eta)}{2}\cdot\Bigl(-\frac{h^3}{6}\Bigr)=-\frac{h^3}{12}f''(\eta).
\qedhere
\]
\end{proof}

O erro depende de $f''$: a regra é exata para retas (grau de exatidão 1). Se $f$ é côncava
para cima ($f''>0$), o trapézio fica acima do gráfico e a regra superestima a integral; se
$f''<0$, subestima.

\subsection{Regra repetida}

Com um único trapézio, $h=b-a$ é grande e o erro também. A solução é dividir $[a,b]$ em $n$
subintervalos iguais, de comprimento $h=\dfrac{b-a}{n}$, com nós $x_i=a+ih$
($i=0,1,\dots,n$), e somar as áreas dos $n$ trapézios. Cada nó interno pertence a dois
trapézios, por isso aparece com peso 2:
\[
\boxed{\;T_n=\frac{h}{2}\Bigl[f(x_0)+2f(x_1)+2f(x_2)+\dots+2f(x_{n-1})+f(x_n)\Bigr]\;}
\]

\begin{teorema}[Erro da regra dos trapézios repetida]\label{teo:errotrapr}
Se $f''$ é contínua em $[a,b]$, existe $\eta\in(a,b)$ tal que
\[
\int_a^b f(x)\,dx-T_n=-\frac{(b-a)\,h^2}{12}\,f''(\eta).
\]
Em particular, se $|f''(x)|\le M_2$ em $[a,b]$, então
$\Bigl|\int_a^b f(x)\,dx-T_n\Bigr|\le\dfrac{(b-a)\,h^2}{12}\,M_2$.
\end{teorema}

\begin{proof}
Somando os erros dos $n$ trapézios (Teorema~\ref{teo:errotrap}), o erro total é
$-\dfrac{h^3}{12}\sum_{i=1}^{n}f''(\eta_i)$. A média
$\frac1n\sum f''(\eta_i)$ está entre o mínimo e o máximo de $f''$; pelo Teorema do Valor
Intermediário, ela é igual a $f''(\eta)$ para algum $\eta$. Assim, o erro é
$-\dfrac{h^3}{12}\,n\,f''(\eta)$, e $nh=b-a$.
\end{proof}

O erro é proporcional a $h^2$: dobrar $n$ (dividir $h$ por 2) divide o erro por
aproximadamente $4$.

\begin{exemplo}[: $\int_0^1 e^{-x^2}\,dx$ pela regra dos trapézios]
Com $n=4$, $h=0{,}25$ e os nós $0$; $0{,}25$; $0{,}5$; $0{,}75$; $1$:
\begin{center}
\renewcommand{\arraystretch}{1.3}
\begin{tabular}{c|ccccc}
$i$ & 0 & 1 & 2 & 3 & 4\\ \hline
$x_i$ & $0$ & $0{,}25$ & $0{,}5$ & $0{,}75$ & $1$\\
$f(x_i)$ & $1$ & $0{,}93941307$ & $0{,}77880078$ & $0{,}56978282$ & $0{,}36787944$\\
peso & 1 & 2 & 2 & 2 & 1
\end{tabular}
\end{center}
\begin{align*}
T_4&=\frac{0{,}25}{2}\bigl[1+2(0{,}93941307+0{,}77880078+0{,}56978282)+0{,}36787944\bigr]\\
&=\frac{0{,}25}{2}\cdot5{,}94387278=0{,}74298410.
\end{align*}
O erro real é $|I-T_4|\approx3{,}84\times10^{-3}$. Para a cota, $f''(x)=(4x^2-2)e^{-x^2}$
e $M_2=\max_{[0,1]}|f''(x)|=2$ (em $x=0$), logo
$\dfrac{(b-a)h^2}{12}M_2=\dfrac{0{,}0625}{12}\cdot2\approx1{,}04\times10^{-2}$, coerente com o
erro observado. Dobrando $n$ sucessivamente:
\begin{center}
\renewcommand{\arraystretch}{1.3}
\begin{tabular}{c|c|c|c}
$n$ & $T_n$ & $|I-T_n|$ & razão entre erros consecutivos\\ \hline
1 & $0{,}6839397206$ & $6{,}29\times10^{-2}$ & \\
2 & $0{,}7313702518$ & $1{,}55\times10^{-2}$ & $4{,}07$\\
4 & $0{,}7429840978$ & $3{,}84\times10^{-3}$ & $4{,}02$\\
8 & $0{,}7458656148$ & $9{,}59\times10^{-4}$ & $4{,}01$\\
16 & $0{,}7465845968$ & $2{,}40\times10^{-4}$ & $4{,}00$
\end{tabular}
\end{center}
A razão se aproxima de $4=2^2$, como previsto.
\end{exemplo}

\begin{figure}[h]
\centering
\begin{tikzpicture}
\begin{axis}[width=10cm,height=6cm,axis lines=middle,xmin=-0.1,xmax=1.2,ymin=0,ymax=1.15,
  xtick={0,0.25,0.5,0.75,1},xticklabels={$0$,$0{,}25$,$0{,}5$,$0{,}75$,$1$},
  ytick={0.5,1},yticklabels={$0{,}5$,$1$},grid=both,grid style={gray!20},xlabel={$x$},ylabel={$y$}]
\addplot[fill=laranja!25,draw=laranja,thick] coordinates
  {(0,0)(0,1)(0.25,0.9394)(0.5,0.7788)(0.75,0.5698)(1,0.3679)(1,0)} -- cycle;
\addplot[laranja,thick] coordinates {(0.25,0)(0.25,0.9394)};
\addplot[laranja,thick] coordinates {(0.5,0)(0.5,0.7788)};
\addplot[laranja,thick] coordinates {(0.75,0)(0.75,0.5698)};
\addplot[domain=-0.05:1.15,samples=100,very thick,azul] {exp(-x^2)};
\addplot[only marks,mark=*,mark size=2pt,vermelho] coordinates
  {(0,1)(0.25,0.9394)(0.5,0.7788)(0.75,0.5698)(1,0.3679)};
\end{axis}
\end{tikzpicture}
\caption{Regra dos trapézios repetida para $\int_0^1 e^{-x^2}\,dx$ com $n=4$.}
\label{fig:trap}
\end{figure}

\begin{atividade}{regra dos trapézios repetida}{\arqTrap}
Abra o aplicativo. À esquerda ficam o gráfico de $f$, os trapézios e os controles; à
direita, o passo $h$, a tabela com os nós, os valores $f(x_i)$, os pesos $c_i$ e os produtos
$c_i\,f(x_i)$, a soma $S$, o valor $T_n=\frac h2S$ e o quadro do erro.
\begin{enumerate}[nosep]
  \item Confira, com $n=4$, todos os números do exemplo acima.
  \item Faça $n=2,4,8,16$ e anote o erro $|I-T_n|$. Calcule as razões entre erros consecutivos.
  \item Marque \emph{Mostrar a área exata}. Em que partes do intervalo os trapézios ficam acima do
        gráfico? Relacione com o sinal de $f''(x)=(4x^2-2)e^{-x^2}$.
  \item Digite $f(x)=3x+1$ e depois $f(x)=x^2$. Para qual das duas a regra é exata? Por quê?
  \item Digite $f(x)=x^2$ em $[0,1]$ e confira que o erro real é \emph{igual} à cota. O que há de
        especial com $f''$ nesse caso?
\end{enumerate}
\end{atividade}

% ==================================================================
\section{Regra 1/3 de Simpson}

\subsection{Regra simples}

Agora interpolamos $f$ por uma parábola em três nós igualmente espaçados
$x_0$, $x_1=x_0+h$, $x_2=x_0+2h$. Pela forma de Lagrange,
$p_2(x)=f(x_0)L_0(x)+f(x_1)L_1(x)+f(x_2)L_2(x)$, e basta integrar as bases. Com a mudança de
variável $x=x_0+sh$ ($dx=h\,ds$, $s\in[0,2]$):
\begin{align*}
\int_{x_0}^{x_2}L_0(x)\,dx&=h\int_0^2\frac{(s-1)(s-2)}{2}\,ds=\frac{h}{3},&
\int_{x_0}^{x_2}L_1(x)\,dx&=h\int_0^2\frac{s(s-2)}{-1}\,ds=\frac{4h}{3},\\
\int_{x_0}^{x_2}L_2(x)\,dx&=h\int_0^2\frac{s(s-1)}{2}\,ds=\frac{h}{3}.
\end{align*}
Portanto
\[
\int_{x_0}^{x_2}f(x)\,dx\approx\frac{h}{3}\bigl[f(x_0)+4f(x_1)+f(x_2)\bigr].
\]
O nome vem do fator $h/3$. Os pesos somam $2h$, que é o comprimento do intervalo, como deve
ser para que a fórmula integre exatamente $f\equiv1$.

\begin{teorema}[Erro da regra 1/3 de Simpson simples]
Se $f^{(4)}$ é contínua em $[x_0,x_2]$, existe $\eta\in(x_0,x_2)$ tal que
\[
\int_{x_0}^{x_2}f(x)\,dx-\frac{h}{3}\bigl[f(x_0)+4f(x_1)+f(x_2)\bigr]=-\frac{h^5}{90}\,f^{(4)}(\eta).
\]
\end{teorema}

\begin{observacao}
Esperaríamos um erro com $f'''$, pois interpolamos por um polinômio de grau 2. Mas a
fórmula é exata também para $x^3$: por simetria, para $f(x)=(x-x_1)^3$ a integral em
$[x_0,x_2]$ é zero e a fórmula também dá $\frac h3[(-h)^3+0+h^3]=0$. Assim, a regra $1/3$
tem grau de exatidão 3, e o erro depende de $f^{(4)}$. Esse ``ganho grátis'' acontece em
todas as fórmulas de Newton-Cotes com número ímpar de nós.
\end{observacao}

\subsection{Regra repetida}

Dividimos $[a,b]$ em um número \textbf{par} $n$ de subintervalos, $h=\dfrac{b-a}{n}$, e
aplicamos a regra simples em cada par $[x_0,x_2]$, $[x_2,x_4]$, \dots, $[x_{n-2},x_n]$. Os
nós de índice ímpar são meios de pares (peso 4); os de índice par interno são extremos de
dois pares (peso 2):
\[
\boxed{\;S_n=\frac{h}{3}\Bigl[f(x_0)+4f(x_1)+2f(x_2)+4f(x_3)+\dots+2f(x_{n-2})+4f(x_{n-1})+f(x_n)\Bigr]\;}
\]

\begin{teorema}[Erro da regra 1/3 de Simpson repetida]
Se $f^{(4)}$ é contínua em $[a,b]$ e $n$ é par, existe $\eta\in(a,b)$ tal que
\[
\int_a^b f(x)\,dx-S_n=-\frac{(b-a)\,h^4}{180}\,f^{(4)}(\eta),
\qquad\text{logo}\qquad
\Bigl|\int_a^b f(x)\,dx-S_n\Bigr|\le\frac{(b-a)\,h^4}{180}\,M_4,
\]
com $M_4=\max_{[a,b]}|f^{(4)}(x)|$.
\end{teorema}

(A demonstração é a mesma do Teorema~\ref{teo:errotrapr}: são $n/2$ parcelas
$-\frac{h^5}{90}f^{(4)}(\eta_i)$ e $\frac n2h^5=\frac{(b-a)h^4}{2}$.)
Agora dobrar $n$ divide o erro por aproximadamente $2^4=16$.

\begin{exemplo}[: $\int_0^1 e^{-x^2}\,dx$ pela regra 1/3 de Simpson]
Com $n=4$ usamos os mesmos cinco valores do exemplo dos trapézios, com pesos
$1,4,2,4,1$:
\begin{align*}
S_4&=\frac{0{,}25}{3}\bigl[1+4(0{,}93941307)+2(0{,}77880078)+4(0{,}56978282)+0{,}36787944\bigr]\\
&=\frac{0{,}25}{3}\cdot8{,}96226456=0{,}74685538.
\end{align*}
O erro real é $|I-S_4|\approx3{,}12\times10^{-5}$, cerca de 120 vezes menor que o dos
trapézios com os \emph{mesmos} cinco valores de $f$. Para a cota,
$f^{(4)}(x)=(16x^4-48x^2+12)e^{-x^2}$ e $M_4=12$ (em $x=0$):
$\dfrac{1\cdot0{,}25^4}{180}\cdot12\approx2{,}60\times10^{-4}$.
\begin{center}
\renewcommand{\arraystretch}{1.3}
\begin{tabular}{c|c|c|c}
$n$ & $S_n$ & $|I-S_n|$ & razão\\ \hline
2 & $0{,}7471804289$ & $3{,}56\times10^{-4}$ & \\
4 & $0{,}7468553798$ & $3{,}12\times10^{-5}$ & $11{,}4$\\
8 & $0{,}7468261205$ & $1{,}99\times10^{-6}$ & $15{,}7$\\
16 & $0{,}7468242574$ & $1{,}25\times10^{-7}$ & $15{,}9$
\end{tabular}
\end{center}
\end{exemplo}

\begin{figure}[h]
\centering
\begin{tikzpicture}
\begin{axis}[width=10cm,height=6cm,axis lines=middle,xmin=-0.15,xmax=2.3,ymin=-0.12,ymax=1.15,
  xtick={0,1,2},xticklabels={$x_0$,$x_1$,$x_2$},ytick={0.5,1},yticklabels={$0{,}5$,$1$},
  grid=both,grid style={gray!20},legend style={at={(0.98,0.98)},anchor=north east,font=\small}]
\addplot[domain=0:2,samples=80,fill=laranja!25,draw=laranja,thick] {1-0.773400*x+0.141279*x^2} \closedcycle;
\addplot[domain=-0.1:2.25,samples=120,very thick,azul] {exp(-x^2)};
\legend{parábola $p_2$,$f(x)=e^{-x^2}$}
\addplot[only marks,mark=*,mark size=2pt,vermelho] coordinates {(0,1)(1,0.3679)(2,0.0183)};
\end{axis}
\end{tikzpicture}
\caption{Regra 1/3 de Simpson simples: a área sob a parábola que passa pelos três nós
substitui a área sob $f$. Aqui, $f(x)=e^{-x^2}$ em $[0,2]$ com $h=1$ (escolhido grande
para que a diferença entre as curvas fique visível).}
\label{fig:simpson}
\end{figure}

\begin{atividade}{regra 1/3 de Simpson repetida}{\arqSimpA}
O aplicativo desenha, em cada par de subintervalos, a região sob a parábola que passa pelos
três nós. À direita aparecem a tabela com os pesos $1,4,2,4,\dots,4,1$, a soma, $S_n$ e o
quadro do erro.
\begin{enumerate}[nosep]
  \item Confira os números do exemplo com $n=4$ e compare com a tabela dos trapézios (os
        valores $f(x_i)$ são os mesmos; só os pesos mudam).
  \item Faça $n=2,4,8,16$ e verifique que as razões entre os erros se aproximam de $16$.
  \item Digite $f(x)=x^3-2x+1$ em $[-1,2]$ com $n=2$. O erro é zero? Explique com o grau de
        exatidão. E com $f(x)=x^4$?
  \item Com $f(x)=e^{-x^2}$ em $[0,2]$ e $n=2$, compare o gráfico com a
        Figura~\ref{fig:simpson}. Aumente $n$ e observe as parábolas se confundirem com $f$.
\end{enumerate}
\end{atividade}

% ==================================================================
\section{Regra 3/8 de Simpson}

Interpolando $f$ por uma cúbica em quatro nós igualmente espaçados $x_0,\dots,x_3$ e
integrando as bases de Lagrange (com $x=x_0+sh$, $s\in[0,3]$), obtemos
\[
\int_{x_0}^{x_3}f(x)\,dx\approx\frac{3h}{8}\bigl[f(x_0)+3f(x_1)+3f(x_2)+f(x_3)\bigr],
\qquad\text{erro}=-\frac{3h^5}{80}\,f^{(4)}(\eta).
\]
O grau de exatidão também é 3: o quarto nó não aumenta o grau, apenas melhora a constante
do erro por trecho. Na forma repetida, com $n$ \textbf{múltiplo de 3}, os nós que separam dois
grupos têm peso 2 e os demais nós internos têm peso 3:
\[
\boxed{\;S^{3/8}_n=\frac{3h}{8}\Bigl[f(x_0)+3f(x_1)+3f(x_2)+2f(x_3)+3f(x_4)+\dots+3f(x_{n-1})+f(x_n)\Bigr]\;}
\]
com erro $-\dfrac{(b-a)\,h^4}{80}\,f^{(4)}(\eta)$.

Comparando as cotas, $\frac{1}{80}>\frac{1}{180}$: com o mesmo $h$, a regra $1/3$ costuma ser mais
precisa. A regra $3/8$ é útil quando o número de subintervalos é ímpar (por exemplo, numa
tabela de dados com $n=7$): aplica-se a $3/8$ em três subintervalos e a $1/3$ nos quatro
restantes.

\begin{exemplo}[: $\int_0^1 e^{-x^2}\,dx$ pela regra 3/8 de Simpson]
Com $n=3$, $h=\frac13$ e $f(0)=1$, $f(\frac13)=0{,}89483932$, $f(\frac23)=0{,}64118039$,
$f(1)=0{,}36787944$:
\begin{align*}
S^{3/8}_3&=\frac{3\cdot\frac13}{8}\bigl[1+3(0{,}89483932)+3(0{,}64118039)+0{,}36787944\bigr]\\
&=\frac18\cdot5{,}97593856=0{,}74699232,
\end{align*}
com erro $|I-S^{3/8}_3|\approx1{,}68\times10^{-4}$. Com $n=6$, $S^{3/8}_6=0{,}7468380575$
e o erro cai para $1{,}39\times10^{-5}$ (razão $\approx12$). Com o mesmo $n=6$, a regra $1/3$
dá erro $\approx 6{,}3\times10^{-6}$ (Exercício~\ref{ex:simpson6}).
\end{exemplo}

\begin{atividade}{regra 3/8 de Simpson repetida}{\arqSimpB}
Aqui cada grupo de três subintervalos é substituído pela cúbica que passa pelos quatro nós.
\begin{enumerate}[nosep]
  \item Confira o exemplo com $n=3$ e depois com $n=6$.
  \item Abra também o aplicativo da regra $1/3$ e use $n=6$ e $n=12$ nos dois. Qual regra
        tem o menor erro? A razão entre os erros é próxima de $\frac{180}{80}=2{,}25$?
  \item Teste $f(x)=x^3$ e $f(x)=x^4$ em $[0,3]$ com $n=3$. Confira que o grau de exatidão é 3.
\end{enumerate}
\end{atividade}

% ==================================================================
\section{Fórmulas de Newton-Cotes: resumo}

As três regras são casos da mesma construção: nós igualmente espaçados e integração do
polinômio interpolador. Na tabela, $h$ é o espaçamento e o erro é o da fórmula simples.

\begin{center}
\renewcommand{\arraystretch}{1.55}
\begin{tabular}{l|c|c|c}
fórmula (simples) & nós & grau de exatidão & erro\\ \hline
trapézios: $\frac h2(f_0+f_1)$ & 2 & 1 & $-\frac{h^3}{12}f''(\eta)$\\
$1/3$ de Simpson: $\frac h3(f_0+4f_1+f_2)$ & 3 & 3 & $-\frac{h^5}{90}f^{(4)}(\eta)$\\
$3/8$ de Simpson: $\frac{3h}{8}(f_0+3f_1+3f_2+f_3)$ & 4 & 3 & $-\frac{3h^5}{80}f^{(4)}(\eta)$\\
Boole: $\frac{2h}{45}(7f_0+32f_1+12f_2+32f_3+7f_4)$ & 5 & 5 & $-\frac{8h^7}{945}f^{(6)}(\eta)$
\end{tabular}
\end{center}

\begin{observacao}
Por que não usar um único polinômio de grau alto em todo o intervalo? Pelo mesmo motivo do
fenômeno de Runge (Figura~\ref{fig:runge}): com muitos nós igualmente espaçados o polinômio
interpolador pode oscilar, e, a partir de 9 nós, as fórmulas de Newton-Cotes passam a ter
pesos negativos, o que amplifica erros de arredondamento. Na prática, usamos fórmulas de
grau baixo \emph{repetidas} em muitos subintervalos.
\end{observacao}

% ==================================================================
\section{Quantos subintervalos usar?}

As cotas de erro permitem escolher $n$ \emph{antes} de calcular, desde que saibamos
limitar a derivada que aparece na fórmula.

\begin{exemplo}[: escolha de $n$ para $\int_1^2\frac{dx}{x}=\ln2$ com erro menor que $10^{-4}$]
Para $f(x)=\frac1x$ temos $f''(x)=\frac{2}{x^3}$ e $f^{(4)}(x)=\frac{24}{x^5}$, ambas
decrescentes em $[1,2]$; logo $M_2=2$ e $M_4=24$.

\textbf{Trapézios:}
$\dfrac{(2-1)h^2}{12}\cdot2\le10^{-4}\iff h^2\le6\times10^{-4}\iff h\le0{,}02449$. Como
$h=\frac1n$, precisamos de $n\ge40{,}8$, isto é, $n=41$ subintervalos (42 valores de $f$).

\textbf{Simpson 1/3:}
$\dfrac{(2-1)h^4}{180}\cdot24\le10^{-4}\iff h^4\le7{,}5\times10^{-4}\iff h\le0{,}1655$, isto é,
$n\ge6{,}04$; como $n$ deve ser par, $n=8$ (9 valores de $f$).

Os erros reais são $|\ln2-T_{41}|\approx3{,}7\times10^{-5}$ e
$|\ln2-S_8|\approx7{,}4\times10^{-6}$. As cotas garantem a tolerância, mas são pessimistas:
com $S_6$ o erro já seria $2{,}3\times10^{-5}$.
\end{exemplo}

Quando não conseguimos limitar a derivada (ou $f$ vem de uma tabela), uma alternativa é
comparar duas aproximações: se $T_n$ e $T_{2n}$ concordam em $k$ casas, temos uma boa ideia
da precisão. A próxima seção transforma essa comparação num método muito mais eficiente.

% ==================================================================
\section{Extrapolação de Richardson e método de Romberg}

\subsection{A ideia de Richardson}

Para $f$ suficientemente derivável, a fórmula de Euler-Maclaurin garante que o erro dos
trapézios tem uma expansão em potências \emph{pares} de $h$:
\[
T(h)=I+c_1h^2+c_2h^4+c_3h^6+\cdots,
\]
com $c_1,c_2,\dots$ que não dependem de $h$. Escrevendo a mesma expressão para $h/2$,
\[
T(h/2)=I+c_1\frac{h^2}{4}+c_2\frac{h^4}{16}+\cdots,
\]
e combinando as duas para eliminar o termo em $h^2$:
\[
\frac{4\,T(h/2)-T(h)}{3}=I-\frac{c_2}{4}h^4+\cdots=T(h/2)+\frac{T(h/2)-T(h)}{3}.
\]
O resultado tem erro $O(h^4)$, muito menor, sem nenhuma nova avaliação de $f$. Isso é a
\emph{extrapolação de Richardson}.

\begin{proposicao}\label{prop:richardson}
Para todo $n$, $\dfrac{4\,T_{2n}-T_n}{3}=S_{2n}$: um passo de Richardson aplicado aos
trapézios reproduz a regra $1/3$ de Simpson.
\end{proposicao}

\begin{proof}
Basta verificar num par de subintervalos $[x_0,x_2]$ de comprimento $2h$:
$T_1=h\,(f_0+f_2)$ e $T_2=\frac h2(f_0+2f_1+f_2)$. Então
$\frac{4T_2-T_1}{3}=\frac{h}{3}\bigl(2f_0+4f_1+2f_2-f_0-f_2\bigr)=\frac h3(f_0+4f_1+f_2)$.
Somando sobre os pares, obtemos $S_{2n}$.
\end{proof}

\subsection{O método de Romberg}

Repetindo a ideia, cada nova coluna elimina mais um termo da expansão ($h^4$, depois $h^6$,
\dots). Organizamos tudo numa tabela triangular.

\begin{tcolorbox}[colback=azul!3,colframe=azul!60,title={Algoritmo: método de Romberg},fonttitle=\bfseries]
\begin{enumerate}[nosep]
  \item Para $k=0,1,\dots,K$: calcule $R_{k,0}=T_{2^k}$ (trapézios com $2^k$ subintervalos).
        Cada $T_{2^k}$ reaproveita todos os valores de $f$ de $T_{2^{k-1}}$:
        $T_{2n}=\frac12T_n+h_{2n}\sum f(\text{novos pontos médios})$.
  \item Para $j=1,\dots,k$:
        \[
        R_{k,j}=R_{k,j-1}+\frac{R_{k,j-1}-R_{k-1,j-1}}{4^j-1}.
        \]
  \item Pare quando $|R_{k,k}-R_{k-1,k-1}|$ for menor que a tolerância. A melhor aproximação
        é $R_{K,K}$.
\end{enumerate}
\end{tcolorbox}

Pela Proposição~\ref{prop:richardson}, a coluna $R_{k,1}$ é a regra $1/3$ de Simpson com $2^k$
subintervalos; a coluna $R_{k,2}$ coincide com a regra de Boole.

\begin{exemplo}[: Romberg para $\int_0^1 e^{-x^2}\,dx$ com $K=4$]
\begin{center}
\renewcommand{\arraystretch}{1.3}
\small
\begin{tabular}{c|c|lllll}
$k$ & $n$ & $R_{k,0}=T_n$ & $R_{k,1}$ & $R_{k,2}$ & $R_{k,3}$ & $R_{k,4}$\\ \hline
0 & 1 & $0{,}6839397206$ & & & & \\
1 & 2 & $0{,}7313702518$ & $0{,}7471804289$ & & & \\
2 & 4 & $0{,}7429840978$ & $0{,}7468553798$ & $0{,}7468337098$ & & \\
3 & 8 & $0{,}7458656148$ & $0{,}7468261205$ & $0{,}7468241699$ & $0{,}7468240185$ & \\
4 & 16 & $0{,}7465845968$ & $0{,}7468242574$ & $0{,}7468241332$ & $0{,}7468241326$ & $0{,}7468241331$
\end{tabular}
\end{center}
Por exemplo, $R_{1,1}=0{,}7313702518+\frac{0{,}7313702518-0{,}6839397206}{3}=0{,}7471804289$.
Com os mesmos 17 valores de $f$ usados por $T_{16}$ (erro $2{,}4\times10^{-4}$), o valor
$R_{4,4}$ tem erro de apenas $2{,}8\times10^{-10}$.
\end{exemplo}

\begin{observacao}
Na linha $k=3$, $R_{3,3}$ (erro $1{,}1\times10^{-7}$) é um pouco pior que $R_{3,2}$ (erro
$3{,}7\times10^{-8}$). A expansão $c_1h^2+c_2h^4+\cdots$ só descreve bem o erro quando $h$ já
é pequeno; com $h=\frac18$ o termo seguinte ainda pesa. Por isso o critério de parada compara
linhas sucessivas da diagonal. Se $f$ não for suave (por exemplo, $\sqrt x$ perto de $0$), a
expansão nem vale, e o ganho do método desaparece (Exercício~\ref{ex:raiz}).
\end{observacao}

\begin{atividade}{método de Romberg}{\arqRomberg}
O controle $K$ escolhe quantas linhas da tabela aparecem; o controle $m$ escolhe quais
trapézios ($2^m$ subintervalos) são desenhados.
\begin{enumerate}[nosep]
  \item Com $K=4$, confira a tabela do exemplo e o erro de $R_{4,4}$.
  \item Compare a coluna $R_{k,1}$ com os valores $S_n$ obtidos no aplicativo da regra $1/3$.
  \item Com $K=5$, quantas avaliações de $f$ foram usadas? Qual é o erro de $R_{5,5}$?
  \item Digite $f(x)=\sqrt{x}$ em $[0,1]$ e compare os erros dos trapézios e de Romberg com os
        de $e^{-x^2}$. O que mudou?
\end{enumerate}
\end{atividade}

% ==================================================================
\section{Quadratura de Gauss-Legendre}

\subsection{A ideia}

Nas fórmulas de Newton-Cotes os nós são fixados de antemão e só os pesos são escolhidos;
com $n$ nós, o grau de exatidão é $n-1$ ou $n$. Se escolhermos também os nós, teremos $2n$
parâmetros livres ($x_1,\dots,x_n$ e $w_1,\dots,w_n$) e poderemos exigir exatidão para
$1,x,\dots,x^{2n-1}$, ou seja, grau de exatidão $2n-1$.

Trabalhamos no intervalo padrão $[-1,1]$. Para $n=2$, queremos
$\int_{-1}^{1}g(t)\,dt\approx w_1g(t_1)+w_2g(t_2)$ exata para $g=1,t,t^2,t^3$:
\[
\begin{cases}
w_1+w_2=2\\
w_1t_1+w_2t_2=0\\
w_1t_1^2+w_2t_2^2=\frac23\\
w_1t_1^3+w_2t_2^3=0
\end{cases}
\]
Por simetria, procuramos $t_2=-t_1$ e $w_1=w_2$. A primeira equação dá $w_1=w_2=1$ e a
terceira, $2t_1^2=\frac23$, ou seja, $t_{1,2}=\mp\frac{1}{\sqrt3}$. Assim,
\[
\int_{-1}^{1}g(t)\,dt\approx g\Bigl(-\tfrac{1}{\sqrt3}\Bigr)+g\Bigl(\tfrac{1}{\sqrt3}\Bigr),
\]
exata para cúbicas com apenas \emph{duas} avaliações (a regra $1/3$ de Simpson precisa de três).

Para $n$ qualquer, os nós são as raízes do \emph{polinômio de Legendre} $P_n(t)$ (por isso
o nome Gauss-Legendre) e os pesos são positivos. A tabela traz os primeiros casos.
\begin{center}
\renewcommand{\arraystretch}{1.4}
\begin{tabular}{c|c|c}
$n$ & nós $t_i$ & pesos $w_i$\\ \hline
1 & $0$ & $2$\\
2 & $\pm\frac{1}{\sqrt3}\approx\pm0{,}5773502692$ & $1$\\
3 & $0$;\quad $\pm\sqrt{\frac35}\approx\pm0{,}7745966692$ & $\frac89$;\quad $\frac59$\\
4 & $\pm0{,}3399810436$;\quad $\pm0{,}8611363116$ & $0{,}6521451549$;\quad $0{,}3478548451$\\
5 & $0$;\quad $\pm0{,}5384693101$;\quad $\pm0{,}9061798459$ & $0{,}5688888889$;\quad $0{,}4786286705$;\quad $0{,}2369268851$
\end{tabular}
\end{center}

\subsection{Mudança de intervalo}

Para integrar em $[a,b]$, usamos a mudança de variável afim
$x=\dfrac{b-a}{2}\,t+\dfrac{a+b}{2}$, que leva $[-1,1]$ em $[a,b]$, com
$dx=\frac{b-a}{2}\,dt$:
\[
\boxed{\;\int_a^b f(x)\,dx\approx G_n=\frac{b-a}{2}\sum_{i=1}^{n}w_i\,f(x_i),\qquad
x_i=\frac{b-a}{2}\,t_i+\frac{a+b}{2}.\;}
\]

\begin{teorema}[Erro da quadratura de Gauss-Legendre]
Se $f^{(2n)}$ é contínua em $[a,b]$, existe $\eta\in(a,b)$ tal que
\[
\int_a^b f(x)\,dx-G_n=\frac{(b-a)^{2n+1}\,(n!)^4}{(2n+1)\,\bigl[(2n)!\bigr]^3}\,f^{(2n)}(\eta).
\]
\end{teorema}

\begin{exemplo}[: $\int_0^1 e^{-x^2}\,dx$ por Gauss-Legendre]
Aqui $\frac{b-a}{2}=\frac{a+b}{2}=\frac12$, logo $x_i=\frac{t_i+1}{2}$.

\textbf{Dois pontos:} $x_{1,2}=\frac12\mp\frac{1}{2\sqrt3}\approx0{,}21132487$ e $0{,}78867513$;
\[
G_2=\frac12\bigl[e^{-0{,}21132487^2}+e^{-0{,}78867513^2}\bigr]
=\frac12\bigl[0{,}95632430+0{,}53686508\bigr]=0{,}74659469,
\]
com erro $2{,}29\times10^{-4}$. Com o mesmo número de avaliações, o trapézio simples tem erro
$6{,}3\times10^{-2}$.

\textbf{Três pontos:} $x_2=\frac12$ e $x_{1,3}=\frac12\mp\frac12\sqrt{0{,}6}\approx0{,}11270167$ e $0{,}88729833$;
\[
G_3=\frac12\Bigl[\tfrac59\,f(0{,}11270167)+\tfrac89\,f(0{,}5)+\tfrac59\,f(0{,}88729833)\Bigr]=0{,}74681458,
\]
com erro $9{,}5\times10^{-6}$, cerca de 37 vezes menor que o de Simpson simples
($S_2$, também com três avaliações). Com cinco pontos, $G_5=0{,}7468241268$, com erro
$6\times10^{-9}$.
\end{exemplo}

\begin{atividade}{quadratura de Gauss-Legendre}{\arqGauss}
O aplicativo mostra os nós $t_i$ e os pesos $w_i$, os nós $x_i$ já transportados para $[a,b]$,
o valor $G_{np}$ e uma comparação com os trapézios e com Simpson. No gráfico, a curva
tracejada é o polinômio que passa pelos pontos $(x_i,f(x_i))$, e a área laranja é a sua
integral.
\begin{enumerate}[nosep]
  \item Confira os valores do exemplo para $np=2$ e $np=3$.
  \item Verifique que a área laranja é igual a $G_{np}$. Por que isso acontece? (Pense na
        fórmula como integral de um polinômio interpolador.)
  \item Com $np=2$, digite $f(x)=x^3$ e depois $f(x)=x^4$ em $[0,1]$. Em qual caso o erro é
        zero? Confira o erro do segundo caso com o teorema acima.
  \item Digite $f(x)=1/(1+25x^2)$ em $[-1,1]$ e aumente $np$ de 1 a 5. O polinômio tracejado
        oscila, mas o erro diminui (com oscilações). Para a quadratura de Gauss, sabe-se que
        $G_n\to I$ para toda $f$ contínua, pois os pesos são positivos; com 20 pontos o erro
        já seria $3{,}6\times10^{-4}$.
\end{enumerate}
\end{atividade}

\begin{observacao}
A quadratura de Gauss não usa os extremos $a$ e $b$. Isso permite integrar funções que não
estão definidas nos extremos, como $\int_0^1\frac{\operatorname{sen}x}{x}\,dx$ ou
$\int_0^1\frac{dx}{\sqrt x}$ (esta última converge devagar, mas a fórmula pode ser aplicada).
Por outro lado, os nós são números irracionais, o que torna o método inadequado para dados
tabelados em pontos igualmente espaçados.
\end{observacao}

% ==================================================================
\section{Comparação dos métodos}

A tabela reúne os resultados do exemplo $\int_0^1 e^{-x^2}\,dx$. O custo de um método de
integração numérica se mede, principalmente, pelo número de avaliações de $f$.

\begin{center}
\renewcommand{\arraystretch}{1.35}
\begin{tabular}{l|c|c|c}
método & avaliações de $f$ & aproximação & erro\\ \hline
trapézios, $n=4$ & 5 & $0{,}7429840978$ & $3{,}8\times10^{-3}$\\
trapézios, $n=16$ & 17 & $0{,}7465845968$ & $2{,}4\times10^{-4}$\\
Simpson $1/3$, $n=4$ & 5 & $0{,}7468553798$ & $3{,}1\times10^{-5}$\\
Simpson $3/8$, $n=6$ & 7 & $0{,}7468380575$ & $1{,}4\times10^{-5}$\\
Simpson $1/3$, $n=16$ & 17 & $0{,}7468242574$ & $1{,}2\times10^{-7}$\\
Romberg, $R_{4,4}$ & 17 & $0{,}7468241331$ & $2{,}8\times10^{-10}$\\
Gauss, 3 pontos & 3 & $0{,}7468145842$ & $9{,}5\times10^{-6}$\\
Gauss, 5 pontos & 5 & $0{,}7468241268$ & $6{,}0\times10^{-9}$
\end{tabular}
\end{center}

\begin{figure}[h]
\centering
\begin{tikzpicture}
\begin{loglogaxis}[width=11cm,height=7cm,xlabel={avaliações de $f$},ylabel={erro $|I-\text{aproximação}|$},
  grid=both,grid style={gray!20},xmin=0.8,xmax=80,ymin=1e-9,ymax=0.2,
  legend style={at={(0.98,0.98)},anchor=north east,font=\small}]
\addplot[thick,cL1,mark=*] coordinates {(2,6.2884e-02)(3,1.5454e-02)(5,3.8400e-03)(9,9.5852e-04)(17,2.3954e-04)(33,5.9878e-05)(65,1.4969e-05)};
\addplot[thick,cL2,mark=square*] coordinates {(3,3.5630e-04)(5,3.1247e-05)(9,1.9877e-06)(17,1.2462e-07)(33,7.7946e-09)};
\addplot[thick,cL3,mark=triangle*] coordinates {(1,3.1977e-02)(2,2.2944e-04)(3,9.5486e-06)(4,3.3532e-07)(5,6.0462e-09)};
\legend{trapézios repetida,Simpson $1/3$ repetida,Gauss-Legendre}
\end{loglogaxis}
\end{tikzpicture}
\caption{Erro em função do número de avaliações de $f$ para $\int_0^1 e^{-x^2}\,dx$, em
escala logarítmica. As retas dos trapézios e de Simpson têm inclinações $-2$ e $-4$
(erros proporcionais a $h^2$ e $h^4$).}
\label{fig:compara}
\end{figure}

\paragraph{Em resumo.}
\begin{itemize}[nosep]
  \item A \textbf{regra dos trapézios} é a mais simples e funciona com qualquer tabela de
        dados, mas converge devagar ($h^2$). É imbatível, porém, para funções periódicas
        integradas num período inteiro, caso em que a expansão de Euler-Maclaurin se anula.
  \item A \textbf{regra 1/3 de Simpson} custa o mesmo e converge com $h^4$; é a escolha padrão
        para dados igualmente espaçados com $n$ par. A \textbf{3/8} completa o caso de $n$ ímpar.
  \item \textbf{Romberg} aproveita as mesmas avaliações dos trapézios e atinge precisão alta
        quando $f$ é suave.
  \item \textbf{Gauss-Legendre} dá a maior precisão por avaliação para funções suaves dadas por
        fórmula, mas exige calcular $f$ em pontos irracionais.
  \item Para funções com singularidades ou mudanças bruscas, todos os métodos perdem
        desempenho; usam-se então métodos \emph{adaptativos}, que refinam a malha só onde é
        necessário (é o que fazem os programas profissionais e o comando \texttt{Integral} do
        GeoGebra).
\end{itemize}

% ==================================================================
\section{Exercícios}

\begin{enumerate}[label=\textbf{\arabic*.},leftmargin=*]
\item Calcule $\int_0^{\pi}\operatorname{sen}x\,dx$ pela regra dos trapézios e pela regra
  $1/3$ de Simpson, ambas com $n=4$. Compare com o valor exato $2$.

\item Verifique que a regra dos trapézios tem grau de exatidão 1 e que a regra $1/3$ de
  Simpson tem grau de exatidão 3, testando-as em $1,x,x^2,x^3,x^4$ no intervalo $[0,2]$ com
  a fórmula simples.

\item Quantos subintervalos são necessários para calcular $\int_0^1 e^{-x^2}\,dx$ com erro
  menor que $10^{-6}$ (a) pela regra dos trapézios; (b) pela regra $1/3$ de Simpson? Use
  $M_2=2$ e $M_4=12$.

\item\label{ex:simpson6} Calcule $S_6$ (regra $1/3$, $n=6$) para $\int_0^1 e^{-x^2}\,dx$ e
  compare com $S^{3/8}_6=0{,}7468380575$.

\item A velocidade de um carro (em m/s) foi registrada a cada 10 segundos:
  \begin{center}
  \begin{tabular}{c|ccccccc}
  $t$ (s) & 0 & 10 & 20 & 30 & 40 & 50 & 60\\ \hline
  $v$ (m/s) & 0 & 8{,}2 & 14{,}6 & 19{,}3 & 22{,}4 & 24{,}1 & 24{,}8
  \end{tabular}
  \end{center}
  Estime a distância percorrida, $\int_0^{60}v(t)\,dt$, pela regra dos trapézios e pelas
  regras $1/3$ e $3/8$ de Simpson.

\item Mostre que, ao integrar as bases de Lagrange de quatro nós igualmente espaçados em
  $[x_0,x_3]$, obtêm-se os pesos $\frac{3h}{8},\frac{9h}{8},\frac{9h}{8},\frac{3h}{8}$ da regra
  $3/8$. (Sugestão: use $x=x_0+sh$, $s\in[0,3]$.)

\item Construa a tabela de Romberg com $K=2$ para $\int_0^1\dfrac{4}{1+x^2}\,dx=\pi$.

\item Use Gauss-Legendre com 2 e com 3 pontos para calcular
  $\int_0^{\pi/2}\operatorname{sen}x\,dx=1$.

\item Deduza a fórmula de Gauss-Legendre com $n=1$ (um nó e um peso) exigindo exatidão para
  $1$ e $t$. Que regra conhecida ela é?

\item \textbf{(GeoGebra)} No aplicativo de Gauss-Legendre, use $np=2$, $f(x)=x^4$ e
  $[a,b]=[0,1]$. Calcule o erro pela fórmula do teorema e compare com o aplicativo.

\item\label{ex:raiz} \textbf{(GeoGebra)} No aplicativo de Romberg, digite $f(x)=\sqrt{x}$ em
  $[0,1]$ (valor exato $\frac23$) e use $K=5$. As razões entre erros consecutivos dos
  trapézios ainda são próximas de 4? O valor $R_{5,5}$ é tão bom quanto no exemplo com
  $e^{-x^2}$? Explique.

\item \textbf{(GeoGebra)} No aplicativo dos trapézios, digite \texttt{sin(2pi x)\^{}2}, isto é,
  $f(x)=\operatorname{sen}^2(2\pi x)$, em $[0,1]$ (valor exato $\frac12$). Compare os resultados
  com $n=2$, $3$, $4$ e $8$ e relacione com a observação sobre funções periódicas da seção de
  comparação.
\end{enumerate}

\subsection*{Respostas}
\begin{enumerate}[label=\textbf{\arabic*.},leftmargin=*,nosep]
\item $T_4\approx1{,}896119$ (erro $0{,}1039$); $S_4\approx2{,}004560$ (erro $0{,}0046$).
\item Trapézios: exata para $1$ e $x$ (dá $2$ e $2$); para $x^2$ dá $4\ne\frac83$.
  Simpson ($h=1$): $\frac13(f(0)+4f(1)+f(2))$ dá $2$, $2$, $\frac83$, $4$, exatos, e para $x^4$
  dá $\frac{20}{3}\ne\frac{32}{5}$.
\item (a) $\frac{h^2}{12}\cdot2\le10^{-6}\Rightarrow h\le0{,}002449\Rightarrow n=409$.
  (b) $\frac{h^4}{180}\cdot12\le10^{-6}\Rightarrow h\le0{,}06223\Rightarrow n\ge16{,}07$, logo $n=18$ (par).
\item $S_6\approx0{,}7468303915$, com erro $\approx6{,}3\times10^{-6}$, cerca de $2{,}2$ vezes menor
  que o da regra $3/8$ com o mesmo $n$ (compare com $\frac{180}{80}=2{,}25$).
\item Trapézios: $1010$ m; Simpson $1/3$: $\approx1017{,}3$ m; Simpson $3/8$: $\approx1017{,}4$ m.
\item Por exemplo, $\int_{x_0}^{x_3}L_0\,dx=h\int_0^3\frac{(s-1)(s-2)(s-3)}{-6}\,ds=\frac{3h}{8}$
  e $\int_{x_0}^{x_3}L_1\,dx=h\int_0^3\frac{s(s-2)(s-3)}{2}\,ds=\frac{9h}{8}$; os outros dois
  seguem por simetria.
\item $R_{0,0}=3$; $R_{1,0}=3{,}1$, $R_{1,1}=3{,}1333333$; $R_{2,0}=3{,}1311765$,
  $R_{2,1}=3{,}1415686$, $R_{2,2}=3{,}1421176$ (erro $5{,}3\times10^{-4}$).
\item $G_2\approx0{,}9984726$ (erro $1{,}5\times10^{-3}$); $G_3\approx1{,}0000081$
  (erro $8{,}1\times10^{-6}$).
\item $w_1=2$ e $w_1t_1=0$, logo $t_1=0$: $\int_a^bf\,dx\approx(b-a)\,f\bigl(\frac{a+b}{2}\bigr)$,
  a regra do ponto médio, com grau de exatidão 1.
\item $G_2=0{,}19\overline{4}$ e o valor exato é $0{,}2$; pelo teorema, o erro é
  $\frac{1\cdot(2!)^4}{5\,(4!)^3}\cdot24=\frac{1}{180}\approx0{,}005556$, exatamente o observado
  (pois $f^{(4)}\equiv24$ é constante).
\item Não: as razões dos trapézios ficam entre $2{,}6$ e $2{,}8$, aproximando-se de $2^{1{,}5}\approx2{,}83$
  (o erro é proporcional a $h^{1{,}5}$), e
  $|I-R_{5,5}|\approx3{,}8\times10^{-4}$, muito pior que no exemplo. Como $f'(x)=\frac{1}{2\sqrt x}$
  não é limitada perto de 0, a expansão de Euler-Maclaurin não vale e a extrapolação não
  elimina os termos certos.
\item Com $n=2$ todos os nós caem nos zeros de $f$ e $T_2=0$ (os nós ``não enxergam'' a
  oscilação); a partir de $n=3$ obtém-se $T_n=\frac12$, exato. Para funções periódicas
  integradas num período completo, os termos da expansão de Euler-Maclaurin se anulam, e a
  regra dos trapézios converge muito rapidamente, desde que haja nós suficientes por período.
\end{enumerate}
